\par\Needspace{8\baselineskip}\bigskip
\pdfbookmark[3]{Exposition, 5}{pallas-II-article5}
\begin{center}5.\end{center}
\noindent\hspace*{\parindent}Il sera à propos d’insérer ici un lemme, que nous n’avons trouvé nulle part énoncé, et dont on peut souvent faire un usage avantageux:
\par\textit{Soient \(Q,R,Q',R'\) quatre points sur la surface d’une sphère, \(\theta\) l’angle entre les grands cercles \(QR,Q'R'\), en considérant les deux branches qui en partant du point de concours vont dans le sens des directions de \(Q\) vers \(R\) et de \(Q'\) vers \(R'\). Cela posé on a}
\[
\begin{aligned}
&\cos QQ'\mathbin{.}\cos RR'-\cos QR'\mathbin{.}\cos RQ'\\
&\qquad=\cos\theta\mathbin{.}\sin QR\mathbin{.}\sin Q'R'.
\end{aligned}
\]
\par On peut substituer à la place de \(\theta\) l’arc \(SS'\), \(S\) et \(S'\) désignant les pôles des grands cercles \(QR,Q'R'\) semblablement posés par rapport à ces grands \SourceLeaf{453}cercles, c’est à dire que le sens du mouvement de \(Q\) vers \(R\) autour de \(S\) soit le même que le sens du mouvement de \(Q'\) vers \(R'\) autour de \(S'\).
\par Appliquant ce théorème aux formules 27—32, on trouve tout de suite
\[
\begin{array}{rl}
33)\quad&Bz-Cy=\cos X\Lambda\mathbin{.}r\BareRoot{mp}\\[4pt]
34)\quad&Cx-Az=\cos Y\Lambda\mathbin{.}r\BareRoot{mp}\\[4pt]
35)\quad&Ay-Bx=\cos Z\Lambda\mathbin{.}r\BareRoot{mp}.
\end{array}
\]
\par Dans le périhélie on a \(\mathrm{d}r=0\), \(r-p=-ae(1-e)\), \(x=r\cos XP=a(1-e)\cos XP\), \(y=a(1-e)\cos YP\), \(z=a(1-e)\cos ZP\). Les formules 21, 22, 23 deviennent donc
\[
\begin{array}{rl}
36)\quad&CG-BH=mpe\cos XP\\[4pt]
37)\quad&AH-CF=mpe\cos YP\\[4pt]
38)\quad&BF-AG=mpe\cos ZP.
\end{array}
\]
\par En les combinant avec les équations 30, 31, 32, notre lemme nous donne
\[
B(BF-AG)-C(AH-CF)=m^{\frac32}p^{\frac32}e\cos X\Pi.
\]
\par Or la même quantité devient \(=F(AA+BB+CC)-A(AF+BG+CH)=mpF\). Ainsi nous avons
\[
\begin{array}{rl}39)\quad&F=e\BareRoot{mp}\mathbin{.}\cos X\Pi\end{array}
\]
et d’une manière semblable
\[
\begin{array}{rl}
40)\quad&G=e\BareRoot{mp}\mathbin{.}\cos Y\Pi\\[4pt]
41)\quad&H=e\BareRoot{mp}\mathbin{.}\cos Z\Pi.
\end{array}
\]
